\section{Latency Analysis}
\label{sec:model}
We model a complete Dispatch or Combine operation as READY, data transfer, and post-barrier. INC starts upload after issuing READY, overlapping data movement with the wait for peer readiness.

\subsection{Communication Latency Model}
\label{sec:delay-model}
Let $T_{\rm ready}$ be the pre-barrier duration, $T_{\rm data}$ the data-transfer phase, and $T_{\rm post}$ the post-barrier and subsequent local epilogue. The data phase includes the sender-side work needed to enter the post-barrier. The baseline executes these stages in sequence:
\begin{equation}
 T_{\rm base}=T_{\rm ready}+T_{\rm data}+T_{\rm post}.
 \label{eq:base}
\end{equation}

Following Swift's separation of processing and network delay~\cite{swift}, consider a reference with simultaneous rank readiness and comparable communication paths:
\begin{equation}
 T_{\rm ready}=T_{\rm local}+T_{\rm signal}.
 \label{eq:ready}
\end{equation}
Here $T_{\rm local}$ includes preparing, issuing, and checking readiness signals. $T_{\rm signal}$ is the network time for the required signal exchange. With concurrent announcements and symmetric paths, this exchange requires one one-way traversal, approximately half a network RTT.

\subsection{Latency Reduction from Early Upload}
\label{sec:readiness-analysis}
In the baseline, ranks finish the entire READY stage before sending DATA. With INC, each rank issues READY and starts uploading. READY and DATA travel toward the network node together; INC holds any early data until the required signals have arrived and destination buffers are writable. \autoref{fig:timing} illustrates this change in timing.

% Schematic timing of endpoint coordination and the proposed \sys{} path.
% Coordinates are schematic, not measurements. Dashed = control; solid = \data{}.
\begin{figure*}[t]
\centering
\begin{tikzpicture}[x=1cm,y=-1cm,font=\footnotesize]
\path[use as bounding box] (0,0) rectangle (17.5,9.3);
\begin{scope}[shift={(0,4.75)}]
\node[paneltitle] at (.1,.15) {(c) Aligned ranks: endpoint exchanges};
\node at (1.1,.65) {Rank 0}; \node at (5.7,.65) {Rank 1};
\draw[lifeline] (1.1,.9)--(1.1,4.2); \draw[lifeline] (5.7,.9)--(5.7,4.2);
\draw[ctrlarrow] (1.1,1.05)--(5.7,1.65);
\draw[ctrlarrow] (5.7,1.05)--(1.1,1.65);
\node[smalllabel] at (3.4,1.08) {\ready{} all-to-all};
\foreach \y in {1.8,2.04,2.28,2.95} {
\draw[dataarrow] (1.1,\y)--(5.7,{\y+.6});
\draw[dataarrow] (5.7,\y)--(1.1,{\y+.6});}
\draw[ctrlarrow] (1.1,3.62)--(5.7,4.17);
\draw[ctrlarrow] (5.7,3.62)--(1.1,4.17);
\node[smalllabel] at (3.4,3.94) {completion exchange};
\draw[<->] (6.1,1.05)--(6.1,1.65) node[midway,right,align=left] {pre\,$\simeq L$};
\draw[<->] (6.1,3.62)--(6.1,4.17) node[midway,right,align=left] {post\,$\simeq L$};
\end{scope}
\begin{scope}[shift={(8.8,4.75)}]
\node[paneltitle] at (.1,.15) {(d) Aligned ranks: \sys{} overlap};
\node at (.8,.65) {Rank 0}; \node at (3.6,.65) {\sys{}}; \node at (6.4,.65) {Rank 1};
\foreach \x in {.8,3.6,6.4} \draw[lifeline] (\x,.9)--(\x,4.2);
\draw[ctrlarrow] (.8,1.05)--(3.6,1.35); \draw[ctrlarrow] (6.4,1.05)--(3.6,1.35);
\fill[incgreen] (3.6,1.35) circle (.045);
\node[smalllabel,anchor=west] at (3.75,1.08) {all \ready{}};
\foreach \y in {1.2,1.44,1.68,2.35} {
\draw[dataarrow] (.8,\y)--(6.4,{\y+.6});
\draw[dataarrow] (6.4,\y)--(.8,{\y+.6});}
\draw[ctrlarrow] (3.6,2.72)--(6.4,3.02);
\draw[ctrlarrow] (3.6,2.72)--(.8,3.02);
\node[smalllabel,align=center] at (3.6,3.44) {\done{} follows final \data{}\\on each ordered output};
\node[anchor=west] at (.4,4.02) {No sender release wait; no new endpoint exchange.};
\end{scope}
\begin{scope}[shift={(0,0)}]
\node[paneltitle] at (.1,.15) {(a) Unequal ranks: endpoint exchanges};
\node at (1.1,.65) {Rank 0}; \node at (5.7,.65) {Rank 1};
\draw[lifeline] (1.1,.9)--(1.1,4.3); \draw[lifeline] (5.7,.9)--(5.7,4.3);
\draw[ctrlarrow] (1.1,1.0)--(5.7,1.6); \draw[ctrlarrow] (5.7,1.55)--(1.1,2.15);
\node[smalllabel] at (3.4,1.02) {Rank 1 ready later};
\foreach \y in {2.2,2.43,2.66,3.03} \draw[dataarrow] (1.1,\y)--(5.7,{\y+.6});
\foreach \y in {1.65,1.88,2.11,2.48} \draw[dataarrow] (5.7,\y)--(1.1,{\y+.6});
\draw[ctrlarrow] (1.1,3.2)--(5.7,3.8);
\draw[ctrlarrow] (5.7,3.72)--(1.1,4.32);
\node[smalllabel,align=center] at (3.4,4.15) {Readiness and completion\\still traverse between endpoints};
\end{scope}
\begin{scope}[shift={(8.8,0)}]
\node[paneltitle] at (.1,.15) {(b) Unequal ranks: bounded staging};
\node at (.8,.65) {Rank 0}; \node at (3.6,.65) {\sys{}}; \node at (6.4,.65) {Rank 1};
\foreach \x in {.8,3.6,6.4} \draw[lifeline] (\x,.9)--(\x,4.3);
\fill[pattern=north east lines,pattern color=black!35] (3.43,1.43) rectangle (3.77,1.87);
\draw[ctrlarrow] (.8,1.0)--(3.6,1.3);
\draw[ctrlarrow] (6.4,1.55)--(3.6,1.85);
\draw[dataarrow] (.8,1.13)--(3.6,1.43);
\draw[dataarrow] (.8,1.36)--(3.6,1.66);
\draw[dataarrow] (.8,1.59)--(3.6,1.89);
\draw[dataarrow] (3.6,1.9)--(6.4,2.2);
\draw[dataarrow] (3.6,2.13)--(6.4,2.43);
\draw[dataarrow] (3.6,2.36)--(6.4,2.66);
\draw[dataarrow] (.8,2.18)--(3.6,2.48)--(6.4,2.95);
\foreach \y in {1.68,1.91,2.14,2.51} \draw[dataarrow] (6.4,\y)--(.8,{\y+.6});
\draw[ctrlarrow] (3.6,2.72)--(6.4,3.02);
\draw[ctrlarrow] (3.6,2.88)--(.8,3.18);
\node[smalllabel,anchor=east] at (3.2,1.68) {buffer};
\node[smalllabel,align=center] at (3.6,3.6) {Wait for \ready{} at the \sys{} node;\\preserve output service and ordering};
\node[anchor=west] at (.4,4.18) {Skew coverage is not an additional fixed saving.};
\end{scope}
\end{tikzpicture}
\caption{Idealized projection of the two-boundary design, not a measured Direct trace. Solid arrows carry \data{}; dashed arrows carry control. The left-hand reference assumes an exposed tail-signal interval of approximately $L$ after data arrival; actual Direct timing follows Equation~\eqref{eq:tailgain}. \sys{} overlaps readiness with upload and appends ordered \done{} at each destination's tail. Early data is buffered under skew. Required local completion and buffer ownership are retained.}
\Description{Four timing panels compare endpoint and network coordination. The top pair shows unequal readiness and buffering; the lower pair shows aligned readiness as an ideal reference. Solid arrows are data and dashed arrows are control.}
\label{fig:timing}
\end{figure*}


To isolate the benefit of overlap, we retain the same local processing cost, data-transfer duration, and post-barrier work in the reference model. When early upload covers the signal-transfer interval, the INC latency is
\begin{equation}
 T_{\rm INC}\approx T_{\rm local}+T_{\rm data}+T_{\rm post}+T_{\rm extra}.
 \label{eq:inc}
\end{equation}
$T_{\rm extra}$ is the additional critical-path delay introduced by INC processing, buffering, and forwarding waits. The ordinary source-to-destination DATA transfer is already included in $T_{\rm data}$; local processing remains accounted for in $T_{\rm local}$.

The reference difference is therefore
\begin{equation}
 \Delta T\approx T_{\rm signal}-T_{\rm extra}.
 \label{eq:gain}
\end{equation}
Early upload thus hides approximately half a network RTT in this reference case, less the added INC delay. The saving comes from overlapping signal transmission with data movement, while the post-barrier remains unchanged.

\subsection{Overlap with Unequal Readiness}
\label{sec:skew}
In practice, differences in computation and scheduling often make ranks ready at different times. \autoref{fig:timing}(C--D) shows how the overlap works in this case: early ranks upload while later ranks prepare, and INC buffers the data. When the last required READY arrives, INC can start forwarding these tokens instead of waiting for the early ranks to receive peer signals and then begin uploading. The mechanism therefore does not require synchronized arrivals.

With sufficient staging and forwarding capacity, this upload head start is preserved under readiness skew and shortens the operation when it advances critical-path data delivery. Output queuing and PFC pauses contribute to $T_{\rm extra}$. The gain depends on the last rank to finish and the available overlap, rather than a fixed fraction of RTT for every arrival pattern.
